\clearpage\label{answers}
{\sans\bfseries\LARGE\color{ink} Answers}\par
\footnotesize\linespread{0.94}\selectfont
Try every puzzle first. Where an answer shows a calculation, check the steps, not just the result.

\subhead{Chapter 1}
\textbf{Multix on the missing gram.} No. A gap of 1~g is the size of the scale's error, so it does not show that anything stopped existing. \textbf{Keep box.} Kept: the weight of each part. Set aside: how the mixture looked.

\textbf{1.1} $98+61+20=179$~g. Off by 1~g.\\
\textbf{1.2} (a) M (b) N (c) M (d) N.\\
\textbf{1.3} $70+50+25=145$, a gap of 5~g. No. Reweigh, check for spills and repeat on one scale.\\
\textbf{1.4} Two scales can disagree by a few grams. Use one scale for every weighing.\\
\textbf{1.5} In a mixture the parts keep their own character, even when dissolved parts cannot be seen. Chapter~3 shows this.\\
\textbf{1.6} $40+25+15=80$.\\
\textbf{1.7} Rice. $99$ is more than half of 179, since half is 89.5.\\
\textbf{1.8} Answers vary.\\
\textbf{1.9} No. The difference is 1~g, the same as the smallest step.\\
\textbf{1.10} ``One grain probably rolled off the table, and the scale may be 1~g off.''

\subhead{Chapter 2}
\textbf{Multix on salt and sand.} No. Salt grains and sand grains are about the same size, so the holes do not tell them apart. \textbf{Keep box.} Kept: the split by size or by steel. Set aside: bits that stuck.

\textbf{2.1} (a) S (b) M (c) X (d) X (e) S.\\
\textbf{2.2} $146+60=206$, a gap of 4~g.\\
\textbf{2.3} $15+75=90$. Yes.\\
\textbf{2.4} A clip that is not magnetic, such as a plated or stainless one. The total matched, but the split was wrong.\\
\textbf{2.5} One group used damp sand and one used dry, so the sand differed. Damp sand clogs a sieve.\\
\textbf{2.6} Weigh the whole 1, shake the sieve 2, weigh what fell through 3, add the parts 4.\\
\textbf{2.7} Sieve: size. Magnet: steel or iron.\\
\textbf{2.8} Either order works, since each tool ignores the other parts. Magnet first keeps clips from catching in the sieve.\\
\textbf{2.9} 35.\\
\textbf{2.10} Answers vary.

\subhead{Chapter 3}
\textbf{Multix on the weights.} The cup weighs the same as before stirring, so it weighs more than 100~g of water plus the cup. The salt's weight stayed. \textbf{Keep box.} Kept: the total weight. Set aside: where the salt sits.

\textbf{3.1} $40+100+7=147$~g.\\
\textbf{3.2} $33-28=5$~g.\\
\textbf{3.3} Put a spoonful of each on a saucer and let it dry. The salty one leaves a crust.\\
\textbf{3.4} ``Five spoonfuls dissolved in 100~g of water and the sixth stayed as grains.''\\
\textbf{3.5} The water temperature and the person who added the salt (and stirred).\\
\textbf{3.6} The empty circle is 100, the water that leaves as vapour.\\
\textbf{3.7} (a) S (b) G (c) B.\\
\textbf{3.8} Sugar dissolves and sand does not. Stir and look for grains, or let it dry and look for a crust.\\
\textbf{3.9} Some crust sticks to the saucer, and the scale is only exact to about 1~g.\\
\textbf{3.10} Answers vary.

\subhead{Chapter 4}
\textbf{Multix on pure water.} No. A filter catches what is too big to pass. Dissolved salt would pass through. \textbf{Keep box.} Kept: the mud's weight. Set aside: anything dissolved.

\textbf{4.1} $185+13=198$, a gap of 2~g.\\
\textbf{4.2} $17-2=15$~g.\\
\textbf{4.3} (a) F (b) E (c) S.\\
\textbf{4.4} Dry a spoonful on a saucer. A crust means something was dissolved.\\
\textbf{4.5} The filters were different papers. The paper towel may have bigger holes. ``Cannot be filtered'' is too big a claim.\\
\textbf{4.6} $150-140-8=2$.\\
\textbf{4.7} 2~g. No. Check the cup and funnel for water and reweigh.\\
\textbf{4.8} (a) S (b) G (c) B.\\
\textbf{4.9} The mud is heavier than the water and settles to the bottom over time.\\
\textbf{4.10} Answers vary.

\subhead{Chapter 5}
\textbf{Multix on filtering first.} The salt is dissolved only after water is added, so a dry mixture cannot be filtered apart. Dissolve first. \textbf{Keep box.} Kept: the order of the steps and the total. Set aside: small losses.

\textbf{5.1} Weigh the whole 1, magnet 2, add water and stir 3, filter 4, then evaporating the salt water 5 and weighing the dry sand 6 (5 and 6 may swap).\\
\textbf{5.2} $80-10=70$~g.\\
\textbf{5.3} $10+49+19=78$, a gap of 2~g.\\
\textbf{5.4} Dissolved salt passes through the filter. Undissolved salt would stay with the sand.\\
\textbf{5.5} Clips: magnet. Sand: filter. Salt: evaporation.\\
\textbf{5.6} Sand stuck to the filter or the cup, or a pinch dropped.\\
\textbf{5.7} Salt crust stuck to the saucer, or salt water left in the cup.\\
\textbf{5.8} No. 9~g is much more than the scale's error. Retrace the steps and weigh after each.\\
\textbf{5.9} Magnet for the clips, add water and dissolve the salt, sieve or pick out the rice, evaporate the salt water.\\
\textbf{5.10} The sand was drier after a day, so less water was left in it.

\subhead{Chapter 6}
\textbf{Multix on the dust.} No. The ledger showed the gap at the sieve, and a pile of sand on the table matched. \textbf{Keep box.} Kept: where the loss happened. Set aside: the cause of the loss.

\textbf{6.1} Step 2. 5~g.\\
\textbf{6.2} $90-15=75$~g.\\
\textbf{6.3} No. A 1~g gap is within the scale's error, so it cannot be traced.\\
\textbf{6.4} Weighing after every step shows the step where a gap appeared.\\
\textbf{6.5} A pile of sand on the table, sand stuck in the sieve, or sand in the bowl.\\
\textbf{6.6} $0+2+1=3$~g.\\
\textbf{6.7} (a) M (b) R (c) M.\\
\textbf{6.8} No. The ledger shows where, not why. You would need to watch or repeat the pour.\\
\textbf{6.9} Different scales can disagree, so the gap could come from the scales.\\
\textbf{6.10} $12+30+53=95$~g.\\
\textbf{6.11} ``5~g of sand went missing at the sieve step and was found on the table.''\\
\textbf{6.12} Weigh after every step, use one scale, or name each gap.\\
\textbf{6.13} Calling the leftover ``salt'' hides any loss. You should weigh the salt.\\
\textbf{6.14 to 6.16} Answers vary. Check that every circle's weight is accounted for.\\
\textbf{Slop check.} (a) You cannot tell: no title, no labels. (b) Many colours, stars, joke text. (c) A title, the parts, and the weights with units. (d) The stars and jokes.

\clearpage\label{notes}
{\sans\bfseries\LARGE\color{ink} Words}\par
\small
\begin{multicols}{2}
\word{dissolve}{to spread out in a liquid until the parts cannot be seen.}
\word{evaporate}{to turn from liquid into vapour and leave.}
\word{filter}{paper with tiny holes that lets liquid pass but catches bits.}
\word{ledger}{a list of weights that adds up to a total.}
\word{mixture}{two or more things put together that keep their own character.}
\word{model}{a simple stand-in that keeps some things and drops others.}
\word{sieve}{a net or screen with holes that lets small bits fall through.}
\word{solution}{a mixture of a liquid with something dissolved in it.}
\end{multicols}

\vspace{0.4em}
{\sans\bfseries\large\color{sky!70!black} Notes for grown-ups}\par
This book teaches that a separation is checked by a ledger of weights: the parts, added up, should match the whole within the scale's error, and weighing after every step shows where a gap appeared. Each chapter has one activity, one worked example with made-up practice numbers, one trust box, one ``what did we keep?'' prompt and a page of puzzles.

\textbf{Materials and safety.} A kitchen scale in grams, magnets, steel paper clips, sand, salt, rice, lentils, soil, filters, cups and a saucer. Never taste an experiment. Keep small magnets and paper clips away from young children and pets. Some clips are plated and not magnetic.

\textbf{Scale error.} A kitchen scale reading to 1~g can be off by about 1~g. The book treats gaps of that size as unexplained and larger gaps as worth tracing.

\textbf{The practice numbers} were made up so the arithmetic is easy to check. Real results will vary with the equipment and the materials.

\subhead{Sources}
\begin{itemize}[leftmargin=1.3em,itemsep=0.2em]
\item T. L. Brown, H. E. LeMay and others. \emph{Chemistry: The Central Science}. Pearson. For mixtures, solutions and methods of separation.
\item J. R. Taylor. \emph{An Introduction to Error Analysis}, 2nd edition. University Science Books, 1997. For measurement error and repeats.
\end{itemize}

\clearpage
\thispagestyle{blank}
\begin{tikzpicture}[remember picture,overlay]
  \fill[sun!35] ($(current page.south west)+(0.4in,0.4in)$) rectangle ($(current page.north east)-(0.4in,0.4in)$);
  \node[anchor=north,align=center,font=\sffamily\bfseries,text=ink,text width=5.4in] at ($(current page.north)+(0,-1.4in)$) {\Large Next in Science\\[6pt]\fontsize{28}{30}\selectfont Sorting the Living};
  \node[align=center,font=\sffamily,text width=4.6in,text=ink] at (current page.center) {A tray holds leaves, shells and buttons. Which belong together? Book 7 builds a yes-or-no key, trades it with a friend, and tests whether it sorts the same way twice.};
  \node at ($(current page.south)+(0,2.6in)$) {\scalebox{3}{\multixbot[ok]}};
  \node[font=\sffamily\bfseries,text=ink] at ($(current page.south)+(0,0.9in)$) {Ages 9 to 11 \quad$\star$\quad Science Book 6};
\end{tikzpicture}
